\subsection{Illustrative Income Counterfactuals}
\label{subsec:income-counterfactuals}

Table~\ref{tab:income-counterfactuals-2025} applies the pooled M2 coefficients
to the 22.56 million workers represented by the 2025 sample, using
Equation~(\ref{eq:income-counterfactual}). Assigning the 12.43 million
informal workers the formal prediction at their current size raises their
mean predicted weekly income by 45.8 percent and total predicted weekly
income across all workers by 14.7 percent. Assigning formal workers to
formal 101+ employment reassigns 4.57 million workers. Their mean predicted
income rises by 29.0 percent, while the total across all workers rises by
7.7 percent. Assigning informal workers to informal 101+ employment reassigns
12.31 million workers, with corresponding increases of 37.7 and 11.9 percent.

\begin{table}[htbp]
\centering\small
\caption{Illustrative income changes under M2, 2025}
\label{tab:income-counterfactuals-2025}
\begin{tabular}{p{0.38\textwidth}rrr}
\toprule
Assignment & \shortstack{Reassigned\\workers\\(millions)} & \shortstack{Reassigned workers\\Mean income\\change (\%)} & \shortstack{All workers\\Total income\\change (\%)} \\
\midrule
Informal to formal, same size & 12.43 & 45.8 & 14.7 \\
Formal to formal 101+ & 4.57 & 29.0 & 7.7 \\
Informal to informal 101+ & 12.31 & 37.7 & 11.9 \\
\midrule
All workers to formal 101+ & 17.00 & 76.6 & 44.9 \\
\bottomrule
\end{tabular}
\par\vspace{0.3em}
\begin{minipage}{0.97\textwidth}\footnotesize
Notes: Equation~(\ref{eq:income-counterfactual}) with pooled M2 coefficients. The sample represents 22.56 million workers. Income is predicted weekly labor income. Hours, employment category, other covariates, coefficients, and fixed effects are unchanged. Workers already at the destination are not counted as reassigned. The common retransformation factor cancels.
\end{minipage}
\end{table}


The joint scenario assigns all workers to formal 101+ employment within
their original status in employment. It reassigns 17.00 million workers,
raising their mean predicted weekly income by 76.6 percent and the total
across all workers by 44.9 percent. Formal workers already in 101+ retain
their baseline predictions. Keeping the population and weights unchanged
means that the percentage change in total income also equals the percentage
change in average income across all workers. The joint increase cannot be
obtained by adding the preceding scenarios, because size and formality
interact and their wage ratios compound.

This joint result is larger than the 25.9 percent obtained under M0's
common-profile restriction. Appendix Table~\ref{tab:status-income-contributions}
shows where the M2 change comes from. Own-account workers contribute
28.8 percentage points and employers 9.1 points to the 44.9 percent total.
Private-sector employees contribute 5.7 points. The large own-account
contribution must be read with GEIH's size definition in mind. Assignment
to own-account work associated with a large organization is not the same
as turning a solo business into a firm with more than 100 employees.

Appendix Table~\ref{tab:income-counterfactual-sensitivity} replaces fitted
baseline income with observed income while applying the same M2 wage ratios.
The joint increase is 45.6 percent under this alternative, compared with
44.9 percent using fitted income. Appendix
Table~\ref{tab:income-counterfactual-support} checks destination-cell support.
Requiring at least 30 pooled observations with the same sector, department,
status in employment, destination size, and formality reduces the joint
increase to 26.1 percent. Requiring 30 observations in 2025 reduces it to
6.3 percent. These screens retain the same coefficients but leave more
workers at their original assignments. They therefore change the extent
of reallocation rather than the estimated gain from an identical move.
\FloatBarrier
